\documentclass[11pt]{article}
\usepackage[margin=0.9in]{geometry}
\usepackage[utf8]{inputenc}
\usepackage[T1]{fontenc}
\usepackage{graphicx}
\usepackage{booktabs}
\usepackage{array}
\usepackage{xcolor}
\usepackage{listings}
\usepackage{hyperref}
\usepackage{parskip}
\definecolor{codebg}{RGB}{245,245,245}
\lstset{basicstyle=\ttfamily\footnotesize, backgroundcolor=\color{codebg},
  breaklines=true, frame=single, columns=fullflexible, keepspaces=true,
  showstringspaces=false}

\title{DATA-260 Homework 5 -- Report}
\author{Uday Patel}
\date{}

\begin{document}
\maketitle

\section*{Section 0 -- Personal Configuration and Domain}
\begin{tabular}{|l|l|}\hline
\textbf{Value} & \textbf{Result} \\\hline
SID4 & 9871 \\\hline
PORT\_BASE & 8871 \\\hline
PREFIX & s9871 \\\hline
SEED & 9871 \\\hline
VERIFY\_SEED & 269871 \\\hline
DOMAIN\_ID & 7 -- Community sports league fixtures \\\hline
\end{tabular}
\vspace{1em}

\textbf{Hardware:} Lenovo Legion Slim 5 16" (Ryzen 5 7640HS, 16GB RAM, RTX 4060 8GB, 512GB SSD). \\
\textbf{Local model used:} qwen3:8b via Ollama (Part 5 agent); sentence-transformers/all-MiniLM-L6-v2 (not used in this homework). \\
\textbf{Repository:} \url{https://github.com/UdayPate/data260-9871} \\
\textbf{Tagged commit:} PLACEHOLDER\_HASH (to be filled in after the final commit
and hw5 tag) \\ \url{https://github.com/UdayPate/data260-9871/tree/PLACEHOLDER_HASH} \\
\textbf{Collaborator access:} confirmed working for Sbnikitha and supriyaselvanganesan.

\section*{Part 1.I -- Database}

Extended the existing \verb|s9871_rel| database (already holding 5{,}001
\verb|fixtures| rows from HW4) rather than creating a new service. The
related entity is \verb|teams|: one team is the home team of many
fixtures. \verb|fixtures| gained a unique code, a defaulted numeric field,
and a foreign key to \verb|teams| with \verb|ON DELETE RESTRICT|.

\begin{lstlisting}
class Team(Base):
    __tablename__ = "teams"
    id = Column(Integer, primary_key=True, autoincrement=True)
    team_name = Column(String(255), nullable=False)             # primary text field
    home_ground = Column(String(255), nullable=False)           # secondary text field
    team_code = Column(String(20), nullable=False, unique=True) # unique field, TM-####
    created_at = Column(DateTime, nullable=False, server_default=func.now())
    updated_at = Column(DateTime, nullable=False, server_default=func.now(),
                         onupdate=func.now())
    fixtures = relationship("Fixture", back_populates="home_team")


class Fixture(Base):
    __tablename__ = "fixtures"
    id = Column(Integer, primary_key=True, autoincrement=True)
    fixture_name = Column(String(255), nullable=False)
    teams_players = Column(String(255), nullable=False)
    fixture_code = Column(String(24), nullable=False, unique=True)  # FX-9871-#####
    spots_available = Column(Integer, nullable=False, server_default="22", default=22)
    home_team_id = Column(Integer, ForeignKey("teams.id", ondelete="RESTRICT"),
                           nullable=False)
    created_at = Column(DateTime, nullable=False, server_default=func.now())
    updated_at = Column(DateTime, nullable=False, server_default=func.now(),
                         onupdate=func.now())
    home_team = relationship("Team", back_populates="fixtures")
\end{lstlisting}

A fixture involves two sides, but the assignment asks for a single foreign
key, so \verb|home_team_id| deliberately models only the home side; the
away side stays in the free-text \verb|teams_players| column.

Because \verb|fixtures| already held 5{,}001 rows, \verb|migrate_hw05.py| adds the
new columns NULL-able, backfills them deterministically from the primary
key, then tightens to NOT NULL and adds the UNIQUE key and foreign key --
\verb|create_all()| cannot ALTER an existing table. Real schema after migration
(no named tool requested a screenshot here, so the real console output
substitutes):

\begin{lstlisting}
CREATE TABLE `teams` (
  `id` int NOT NULL AUTO_INCREMENT,
  `team_name` varchar(255) NOT NULL,
  `home_ground` varchar(255) NOT NULL,
  `team_code` varchar(20) NOT NULL,
  `created_at` datetime NOT NULL DEFAULT (now()),
  `updated_at` datetime NOT NULL DEFAULT (now()),
  PRIMARY KEY (`id`),
  UNIQUE KEY `team_code` (`team_code`)
) ENGINE=InnoDB AUTO_INCREMENT=13 DEFAULT CHARSET=utf8mb4 COLLATE=utf8mb4_0900_ai_ci;

CREATE TABLE `fixtures` (
  `id` int NOT NULL AUTO_INCREMENT,
  `fixture_name` varchar(255) NOT NULL,
  `teams_players` varchar(255) NOT NULL,
  `fixture_code` varchar(24) NOT NULL,
  `spots_available` int NOT NULL DEFAULT '22',
  `home_team_id` int NOT NULL,
  `created_at` datetime NOT NULL DEFAULT CURRENT_TIMESTAMP,
  `updated_at` datetime NOT NULL DEFAULT CURRENT_TIMESTAMP ON UPDATE CURRENT_TIMESTAMP,
  PRIMARY KEY (`id`),
  UNIQUE KEY `uq_fixtures_fixture_code` (`fixture_code`),
  KEY `idx_fixture_name` (`fixture_name`),
  KEY `fk_fixtures_home_team` (`home_team_id`),
  CONSTRAINT `fk_fixtures_home_team` FOREIGN KEY (`home_team_id`)
    REFERENCES `teams` (`id`) ON DELETE RESTRICT
) ENGINE=InnoDB AUTO_INCREMENT=5005 DEFAULT CHARSET=utf8mb4 COLLATE=utf8mb4_0900_ai_ci;

row count teams = 12
row count fixtures = 5001
\end{lstlisting}

\section*{Part 1.II -- API Implementation}

Full CRUD for both entities, pagination on the team listing, and a
relationship-query endpoint (shown in the code below). Deleting a
team that still has fixtures is refused at the API (409) and independently
at the database level (the foreign key above):

\begin{lstlisting}
@router.delete("/teams/{team_id}")
def delete_team(team_id: int, user: User = Depends(get_current_user),
                 db: DBSession = Depends(get_db)):
    team = db.query(Team).filter(Team.id == team_id).first()
    if not team:
        raise HTTPException(status_code=404, detail=f"Team {team_id} not found")
    dependents = db.query(Fixture).filter(Fixture.home_team_id == team_id).count()
    if dependents:
        raise HTTPException(status_code=409, detail=(
            f"Team {team_id} still has {dependents} associated fixture(s) "
            "and cannot be deleted"))
    db.delete(team)
    db.commit()
    return {"success": True, "deleted_id": team_id}


@router.get("/teams/{team_id}/fixtures")
def list_fixtures_for_team(team_id: int, user: User = Depends(get_current_user),
                            db: DBSession = Depends(get_db)):
    team = db.query(Team).filter(Team.id == team_id).first()
    if not team:
        raise HTTPException(status_code=404, detail=f"Team {team_id} not found")
    rows = db.query(Fixture).filter(Fixture.home_team_id == team_id)\
             .order_by(Fixture.id).all()
    return {"team": team_to_dict(team), "fixture_count": len(rows),
            "fixtures": [fixture_to_dict(f) for f in rows]}
\end{lstlisting}

Tested against the real server and the real \verb|s9871_rel| database with
a 15-request Postman collection, run twice consecutively (15/15 both
times), fully self-cleaning.

\includegraphics[width=0.8\textwidth]{screenshots/p1_team_create.png} \\
\small \textit{Figure 1: Postman -- create team (201).}

\includegraphics[width=0.8\textwidth]{screenshots/p1_team_list_paginated.png} \\
\small \textit{Figure 2: Postman -- list teams, paginated (200).}

\includegraphics[width=0.8\textwidth]{screenshots/p1_team_get_one.png} \\
\small \textit{Figure 3: Postman -- get one team (200).}

\includegraphics[width=0.8\textwidth]{screenshots/p1_team_update.png} \\
\small \textit{Figure 4: Postman -- update team (200).}

\includegraphics[width=0.8\textwidth]{screenshots/p1_fixture_create.png} \\
\small \textit{Figure 5: Postman -- create fixture, spots\_available defaulted to 22 (201).}

\includegraphics[width=0.8\textwidth]{screenshots/p1_fixture_list.png} \\
\small \textit{Figure 6: Postman -- list fixtures (200).}

\includegraphics[width=0.8\textwidth]{screenshots/p1_fixture_get_one.png} \\
\small \textit{Figure 7: Postman -- get one fixture (200).}

\includegraphics[width=0.8\textwidth]{screenshots/p1_fixture_update.png} \\
\small \textit{Figure 8: Postman -- update fixture, spots 22 to 14 (200).}

\includegraphics[width=0.8\textwidth]{screenshots/p1_relationship_query.png} \\
\small \textit{Figure 9: Postman -- relationship query, fixtures for one team (200).}

\includegraphics[width=0.8\textwidth]{screenshots/p1_error_404_not_found.png} \\
\small \textit{Figure 10: Postman -- resource-not-found error (404).}

\includegraphics[width=0.8\textwidth]{screenshots/p1_error_422_validation.png} \\
\small \textit{Figure 11: Postman -- validation error on the unique-field format (422).}

\includegraphics[width=0.8\textwidth]{screenshots/p1_error_409_constraint.png} \\
\small \textit{Figure 12: Postman -- constraint violation, delete refused while fixtures exist (409).}

\includegraphics[width=0.8\textwidth]{screenshots/p1_fixture_delete.png} \\
\small \textit{Figure 13: Postman -- delete fixture (200).}

\includegraphics[width=0.8\textwidth]{screenshots/p1_team_delete.png} \\
\small \textit{Figure 14: Postman -- delete team, now allowed (200).}

\section*{Part 1.III -- Redux Client}

Redux Toolkit store with a \verb|fixtures| slice (the primary domain
entity) and a small read-only \verb|teams| slice for the home-team
dropdown. Four async thunks call the FastAPI endpoints via axios.

\begin{lstlisting}
export const store = configureStore({
  reducer: { fixtures: fixturesReducer, teams: teamsReducer },
});

export const createFixture = createAsyncThunk("fixtures/create",
  async (fixture, { rejectWithValue }) => {
    try {
      const response = await http.post("/fixtures", fixture);
      return response.data;
    } catch (err) { return rejectWithValue(errorMessage(err)); }
  });

export const updateFixture = createAsyncThunk("fixtures/update",
  async ({ id, changes }, { rejectWithValue }) => {
    try {
      const response = await http.put(`/fixtures/${id}`, changes);
      return response.data;
    } catch (err) { return rejectWithValue(errorMessage(err)); }
  });

export const deleteFixture = createAsyncThunk("fixtures/delete",
  async (id, { rejectWithValue }) => {
    try {
      await http.delete(`/fixtures/${id}`);
      return Number(id);  // reducer needs the id to drop the row
    } catch (err) { return rejectWithValue(errorMessage(err)); }
  });
\end{lstlisting}

Home reads everything from the store via \verb|useSelector|, so the table
updates itself after every create/update/delete with no manual refetch:

\begin{lstlisting}
export default function Home({ user }) {
  const dispatch = useDispatch();
  const fixtures = useSelector(selectFixtures);
  const status = useSelector(selectFixturesStatus);
  const lastTouchedId = useSelector(selectLastTouchedId);
  const teams = useSelector(selectTeams);
  // ... pagination (25/page) and per-row Edit/Delete rendered below
\end{lstlisting}

Verified with a real headless-browser test (Playwright): login, pagination,
create + direct MySQL check, a forced 409 and 422 shown on the form, update
by id with \verb|updated_at| advancing, and per-row delete confirmed gone
from both the table and MySQL -- 22/22, twice consecutively.

Per CLAUDE.md's established convention, ``code snippet and UI output
together, in the same frame'' is satisfied here in the report layout: the
relevant snippet directly above each screenshot, one caption below both.

\includegraphics[width=0.8\textwidth]{screenshots/p1_redux_home.png} \\
\small \textit{Figure 15: Home screen -- reads the fixture list from Redux state (code above).}

\includegraphics[width=0.8\textwidth]{screenshots/p1_redux_create.png} \\
\small \textit{Figure 16: Create screen -- dispatches createFixture (code above).}

\includegraphics[width=0.8\textwidth]{screenshots/p1_redux_update.png} \\
\small \textit{Figure 17: Update screen -- select by id, dispatches updateFixture (code above).}

\includegraphics[width=0.8\textwidth]{screenshots/p1_redux_delete.png} \\
\small \textit{Figure 18: Delete -- per-row delete button dispatches deleteFixture (code above).}

\section*{Part 2.A -- TheMealDB Tutorial MCP Server}

A local MCP server named \verb|meals| (\verb|code/meals_server.py|), four
tools over STDIO against the real TheMealDB API. mcp 2.2.0 renamed
FastMCP to MCPServer; the API is otherwise the same.

\begin{lstlisting}
def meal_details(id: str | int) -> dict | str:
    """Look up one meal by its TheMealDB id and return the full recipe."""
    meal_id = str(id).strip()
    if not meal_id.isdigit():
        raise ToolError(f"id must be numeric, got {id!r}")
    data = _get("lookup.php", {"i": meal_id})
    meals = data.get("meals")
    if not meals:
        return f"no meal found with id {meal_id}"
    return _full(meals[0])
\end{lstlisting}

Errors are raised as \verb|ToolError| specifically -- in mcp 2.2.0 that is
the only exception type whose message reaches the MCP client; others are
masked as a generic string, which would hide the real reason in the
Inspector.

\includegraphics[width=0.8\textwidth]{screenshots/p2a_search_meals_by_name.png} \\
\small \textit{Figure 19: MCP Inspector -- search\_meals\_by\_name(Arrabiata, 5) -- Spicy Arrabiata Penne.}

\includegraphics[width=0.8\textwidth]{screenshots/p2a_meals_by_ingredient.png} \\
\small \textit{Figure 20: MCP Inspector -- meals\_by\_ingredient.}

\includegraphics[width=0.8\textwidth]{screenshots/p2a_meal_details.png} \\
\small \textit{Figure 21: MCP Inspector -- meal\_details(52771), 8 ingredients.}

\includegraphics[width=0.8\textwidth]{screenshots/p2a_random_meal.png} \\
\small \textit{Figure 22: MCP Inspector -- random\_meal, same shape as meal\_details.}

\section*{Part 2.B -- Domain MCP Server}

A second server, \verb|domain|, over the real \verb|s9871_rel| data: exactly
three tools, search / detail lookup / aggregate, all returning one shared
envelope.

\begin{lstlisting}
def result(data: Any = None, error: str | None = None) -> dict:
    """The one envelope shape every domain tool returns."""
    return {"ok": error is None, "data": data, "error": error}

@mcp.tool()
def team_fixture_counts(min_fixtures: int = 0) -> dict:
    """Aggregate: number of fixtures hosted by each team."""
    if min_fixtures < 0:
        return result(error=f"min_fixtures must be >= 0, got {min_fixtures}")
    db = db_session_basede26()
    try:
        teams = db.query(Team).order_by(Team.id).all()
        rows = db.query(Fixture.home_team_id, func.count(Fixture.id))\
                 .group_by(Fixture.home_team_id).all()
        counts = dict(rows)
        data = [{"team_id": t.id, "team_code": t.team_code, "team_name": t.team_name,
                 "fixture_count": counts.get(t.id, 0)}
                for t in teams if counts.get(t.id, 0) >= min_fixtures]
        return result(data=data)
    finally:
        db.close()
\end{lstlisting}

Validated in MCP Inspector with one successful and one intentionally
invalid call per tool, against the live database:

\includegraphics[width=0.8\textwidth]{screenshots/Search_fixtures_valid.png} \\
\small \textit{Figure 23: search\_fixtures("lions") -- ok: true, 5 real matches.}

\includegraphics[width=0.8\textwidth]{screenshots/Search_fixtures_invalid.png} \\
\small \textit{Figure 24: search\_fixtures("") -- ok: false, "query must not be empty".}

\includegraphics[width=0.8\textwidth]{screenshots/fixture_details_valid.png} \\
\small \textit{Figure 25: fixture\_details(4) -- ok: true, nested home\_team.}

\includegraphics[width=0.8\textwidth]{screenshots/Fixture_details_invalid.png} \\
\small \textit{Figure 26: fixture\_details(999999) -- ok: false, "fixture 999999 not found".}

\includegraphics[width=0.8\textwidth]{screenshots/team_fixtures_valid.png} \\
\small \textit{Figure 27: team\_fixture\_counts(min\_fixtures=100) -- ok: true, real per-team counts.}

\includegraphics[width=0.8\textwidth]{screenshots/team_fixures_invalid.png} \\
\small \textit{Figure 28: team\_fixture\_counts(min\_fixtures=-2) -- ok: false, "min\_fixtures must be >= 0, got -2".}

\section*{Part 3 -- Tool Contracts Under Stress}

\textbf{3.18 -- expected schema, rejected input, returned error} (same three
tools and rejected calls as Part 2.B above):

{\small\begin{tabular}{|p{3.4cm}|p{4.0cm}|p{3.0cm}|p{3.8cm}|}\hline
\textbf{Tool} & \textbf{Expected schema} & \textbf{Rejected input} & \textbf{Returned error} \\\hline
search\_fixtures & \{query: string (req.), limit: int=10\} & \{"query": ""\} & "query must not be empty" \\\hline
fixture\_details & \{id: integer (req.)\} & \{"id": 999999\} & "fixture 999999 not found" \\\hline
team\_fixture\_counts & \{min\_fixtures: int=0\} & \{"min\_fixtures": -2\} & "min\_fixtures must be >= 0, got -2" \\\hline
\end{tabular}}
\vspace{0.5em}

\textbf{3.19 -- retry/timeout policy}, each attempt in a worker thread so
the timeout is a real enforced cutoff, not just a latency measurement:

\begin{lstlisting}
def call_with_retry(operation, *, max_retries=3, base_delay_s=0.05,
                     max_delay_s=1.0, timeout_s=2.0) -> dict:
    start = time.monotonic()
    last_err = None
    for attempt in range(max_retries + 1):
        future = _executor.submit(operation)
        try:
            data = future.result(timeout=timeout_s)
            return {"ok": True, "data": data, "error": None, "attempts": attempt + 1,
                    "latency_ms": (time.monotonic() - start) * 1000}
        except FutureTimeoutError:
            last_err = f"timed out after {timeout_s}s"
        except Exception as exc:
            last_err = str(exc)
        if attempt < max_retries:
            time.sleep(min(base_delay_s * (2**attempt), max_delay_s))
    return {"ok": False, "data": None, "error": last_err, "attempts": max_retries + 1,
            "latency_ms": (time.monotonic() - start) * 1000}
\end{lstlisting}

Demonstrated, real output:

\begin{lstlisting}
1. success on first attempt: {"ok": true, "data": {"status": "ok"}, "error": null,
   "attempts": 1, "latency_ms": 0.0}
2. failure then success on retry: {"ok": true, "data": {"status": "ok",
   "succeeded_on_attempt": 2}, "error": null, "attempts": 2, "latency_ms": 47.0}
3. failure after all retries: {"ok": false, "data": null, "error": "operation always
   fails", "attempts": 4, "latency_ms": 156.0}
\end{lstlisting}

\textbf{3.20 -- VERIFY\_SEED fault injection}, 150 calls (50 per rate),
VERIFY\_SEED = 269871, raw data in \verb|reports/hw05/raw/|:

\begin{tabular}{|l|l|l|l|}\hline
\textbf{Injected failure rate} & \textbf{Success rate} & \textbf{Mean latency (ms)} & \textbf{p99 latency (ms)} \\\hline
0\%  & 100\% & 13.1 & 31.0 \\\hline
20\% & 100\% & 28.4 & 187.0 \\\hline
50\% & 94\%  & 80.0 & 428.4 \\\hline
\end{tabular}
\vspace{0.5em}

With max\_retries=3 (4 attempts total), the policy fully absorbs 0\% and
20\% injected failure -- failing 4 attempts in a row at 20\% has
probability $0.2^4 \approx 0.16\%$. At 50\%, failing all 4 has probability
$0.5^4 = 6.25\%$, matching the observed 6\% failure almost exactly. p99
grows much faster than the mean because it is dominated by the calls that
need several retries with exponential backoff.

\textbf{Interactive vs.\ batch:} suitable for an interactive assistant up
to about 20\% failure -- every call still succeeds and p99 stays under
200ms. At 50\%, a p99 of 428ms feels sluggish and 1 in 17 requests still
fails outright. For batch processing, increase max\_retries (5--6) and
raise max\_delay\_s (5--10s) so a transient outage gets more chances to
clear -- worse worst-case latency traded for a higher overall success rate,
the right trade when nothing is blocking an interactive user.

\section*{Part 4 -- One Safe Tool Entry Point and Offline Tests}

\verb|execute_tool(name, inputs)| is the single function Part 5's agent
calls. It dispatches through the retry policy above and always returns a
JSON string in the \{ok, data, error\} shape -- never raises, even on an
unknown tool name or a missing argument.

\begin{lstlisting}
def execute_tool(name: str, inputs: dict, tools: dict | None = None) -> str:
    registry = tools or TOOLS
    fn = registry.get(name)
    if fn is None:
        return json.dumps({"ok": False, "data": None,
                            "error": f"unknown tool: {name!r} (known: {sorted(registry)})"})

    if name == "search_fixtures":            # Part 5.I safety rule
        query = (inputs.get("query") or "").strip()
        if 0 < len(query) < 3:
            return json.dumps({"ok": False, "data": None,
                "error": "query too broad: must be at least 3 characters"})

    def call():
        return fn(**inputs)

    outcome = call_with_retry(call, max_retries=2, base_delay_s=0.05, timeout_s=5.0)
    envelope = outcome["data"] if outcome["ok"] else \
        {"ok": False, "data": None, "error": outcome["error"]}
    return json.dumps(envelope)
\end{lstlisting}

A Python test runner using \verb|assert|-backed checks, 7 tests covering
one valid and one invalid input per domain tool plus \verb|execute_tool|'s
own unknown-tool-name dispatch, via dependency injection (fake tool
callables, no MySQL, no LLM, fully offline and repeatable):

\includegraphics[width=0.8\textwidth]{screenshots/p4_execute_tool_tests.png} \\
\small \textit{Figure 29: terminal -- Part 4 offline test runner, 7/7, PASS/FAIL per test and final summary.}

\section*{Part 5 -- Agent Loop, Safety Rule, and Reflection}

\textbf{I. Safety rule} -- a \verb|search_fixtures| query under 3
characters is rejected inside \verb|execute_tool|, before the domain tool
or any DB call runs:

\begin{lstlisting}
allowed: query="Lions" -> {"ok": true, "data": [...2 matches...], "error": null}
blocked: query="a"     -> {"ok": false, "data": null,
                            "error": "query too broad: must be at least 3 characters"}
\end{lstlisting}

\textbf{II. Agent loop} -- \verb|run_agent(user_input)| asks qwen3:8b
(through the existing \verb|ModelClient| adapter from HW1/2) each turn to
either call one of the three tools or give a final answer, in a strict
JSON protocol; a turn counter enforces \verb|max_steps|; every step and the
final stop reason are logged to \verb|agent_runs.jsonl|.

\begin{lstlisting}
class MockModel:
    """Always asks for another tool call - guarantees max_steps is hit."""
    def complete(self, messages, tools=None):
        return json.dumps({"action": "tool", "tool": "team_fixture_counts", "inputs": {}})

while step < max_steps:
    step += 1
    raw = model.complete(messages)
    parsed = _extract_json(raw)   # tolerant of <think> blocks / stray prose
    if parsed.get("action") == "final":
        final_answer = parsed.get("answer", ""); stop_reason = "final_answer"; break
    if parsed.get("action") == "tool":
        result = json.loads(execute_tool(parsed["tool"], parsed.get("inputs") or {},
                                          tools=tools))
        tool_call_count += 1
        log({"step": step, "tool": parsed["tool"], "result": result})
        if not result.get("ok") and "query too broad" in (result.get("error") or ""):
            stop_reason = "safety_rule_block"; break
        messages.append({"role": "user",
                          "content": f"Tool result: {json.dumps(result)}"})
        continue
# loop exits with stop_reason = "max_steps_ceiling" if nothing else set it
\end{lstlisting}

\textbf{III. Additional offline tests} -- two tests added to the existing
Part 4 suite (9 total), both fully offline (fake tools, \verb|MockModel|,
no API key/live model/database):

\begin{lstlisting}
PASS  search_fixtures: valid input -> ok=true
PASS  search_fixtures: invalid (empty query) -> ok=false, documented error
PASS  fixture_details: valid input -> ok=true
PASS  fixture_details: invalid (unknown id) -> ok=false, documented error
PASS  team_fixture_counts: valid input -> ok=true
PASS  team_fixture_counts: invalid (negative min_fixtures) -> ok=false, documented error
PASS  execute_tool: unknown tool name -> ok=false, no crash
PASS  execute_tool: safety rule blocks a too-broad query, no crash
PASS  run_agent + MockModel: stops at max_steps, offline

9/9 tests passed
\end{lstlisting}

\textbf{IV. Metrics} -- 4 real scenarios against local qwen3:8b:

\begin{tabular}{|p{5.3cm}|l|l|l|}\hline
\textbf{Scenario} & \textbf{Stop reason} & \textbf{Steps} & \textbf{Tool calls} \\\hline
A -- detail lookup & final\_answer & 2 & 1 \\\hline
B -- aggregate + follow-up reasoning & final\_answer & 2 & 1 \\\hline
C -- asked to run a too-short search query & final\_answer & 1 & 0 \\\hline
D -- forced ceiling (max\_steps=1) & max\_steps\_ceiling & 1 & 1 \\\hline
\end{tabular}
\vspace{0.5em}

Honestly reported: Scenario A answered correctly (fixture 4's real code and
spots). Scenario B's tool call succeeded with real data, but the model's
final answer was nonsensical -- it claimed the conversation was ``stuck in
a loop'' rather than reporting the aggregate result; this is a genuine
model failure, not a bug in \verb|execute_tool|/\verb|run_agent|, and is
reported as-is rather than re-run for a cleaner result. Scenario C is
notable: the model read the system prompt's stated 3-character rule and
declined the tool call itself (\verb|tool_call_count=0|), so the live agent
never actually exercised the safety-rule code path -- that path is
demonstrated directly above instead. Scenario D is a clean, deliberately
forced \verb|max_steps_ceiling|.

\textbf{V. Reflection} (REFLECTION.md, run\_id \verb|45743a9ea8d1|, scenario
D):

On turn 1, \verb|run_agent| sent the system prompt and the question ``Look
up fixture id 10 and summarize it for me.'' to qwen3:8b. The model replied
\verb|{"action": "tool", "tool": "fixture_details", "inputs": {...}}|,
correctly recognizing the question needed a tool call, not a guess.
\verb|execute_tool| ran \verb|fixture_details(id=10)| through the retry
policy, succeeding on the first attempt and returning the real row:
``Bears vs Panthers (Volleyball \#8)'', code FX-9871-00010, 22 spots
available, home team Santa Clara Bears. That envelope was logged as the
step record and appended to the conversation for the model's next turn.

That is where the run stopped. \verb|max_steps| was deliberately set to 1,
so the loop condition was already false after the one turn that produced
the tool call -- no second turn remained for the model to read the result
and write a summary. \verb|run_agent| exits, finds no final answer and no
safety-rule block, and falls back to \verb|stop_reason = "max_steps_ceiling"|.
The final log record shows exactly this: \verb|tool_call_count: 1|,
\verb|steps: 1|, \verb|final_answer: null| -- the agent had the right data
but ran out of turns before reporting it.

This illustrates why \verb|max_steps| matters even when nothing goes
wrong: a correctly-behaving agent can still hit a step budget too low for
the task. This question needed at least two turns -- call the tool, then
read the result back -- and \verb|max_steps=1| wasn't enough by design. The
harness stopped cleanly rather than looping forever or guessing without
the data it had just fetched.

\end{document}
